% ============================================================
% PHẦN 7 — Slide 61-67 (Quan sát bổ sung, Tổng kết, Q&A)
% ============================================================

% --- Slide 61 ---
\begin{frame}[fragile]{Quan sát 1: \texttt{final\_prune\_structgs} không hề chạy trong pipeline hiện tại}
\begin{itemize}
    \item Đính chính: \texttt{final\_prune\_structgs} \textbf{không} bị khoá bởi cửa sổ iteration $15000$--$30000$ -- không có đoạn code nào áp cửa sổ đó cho hàm này trong \texttt{SADGS/}. (Số $15000$ thực ra thuộc một hàm khác, xem box bên dưới.)
    \item Sự thật trong mã nguồn (\texttt{SADGS/train.py}, quanh dòng 412--419): lệnh gọi \texttt{gaussians.final\_prune\_structgs(...)} \textbf{đã bị comment-out}, đứng cạnh một prune đơn giản đang chạy thật:
\begin{verbatim}
if iteration in prune_iterations:
    prune_mask = (gaussians.get_opacity < 0.1).squeeze()
    gaussians.prune_points(prune_mask)
    # gaussians.final_prune_structgs(min_opacity=0.1,
    #                                 pruning_score=pruning_score)
\end{verbatim}
    \item[$\Rightarrow$] Cơ chế loại floater đa-view (\texttt{final\_prune\_structgs}) là \textbf{dead code} trong pipeline mặc định hiện tại -- không chạy dù ngân sách iteration là bao nhiêu.
    \item Prune thực sự đang chạy chỉ là ngưỡng opacity đơn giản, tại \texttt{prune\_iterations} (mặc định \texttt{[4000, 8000]}). Với \textbf{iterations=7000}, mốc \textbf{8000 không bao giờ tới} $\Rightarrow$ scene chỉ được prune-opacity đúng \textbf{1 lần} (ở iteration 4000) thay vì 2.
    \item Đây là lời giải thích khả dĩ cho: kích thước file lớn bất thường (\textbf{248 B/Gaussian}, chưa nén) và LPIPS "chững" quanh \textbf{0.31}.
    \item \alert{Hệ quả quan trọng:} mọi so sánh trực tiếp giữa số liệu 7k-iteration này với kết quả công bố gốc của SADGS/3DGS (chạy đủ 30k iteration, và giả định \texttt{final\_prune\_structgs} có được kích hoạt) đều \textbf{không hợp lệ} — khác ngân sách huấn luyện, khác điểm dừng pipeline, khác cả việc floater-pruning cuối có chạy hay không.
\end{itemize}
\vspace{1mm}
{\scriptsize \textit{Ghi chú:} số $15000$ trong \texttt{SADGS/scene/gaussian\_model.py} (\texttt{optimizer\_step}) là ngưỡng đổi tần suất gọi \texttt{optimizer.step()} (mỗi iteration $\to$ mỗi 32 $\to$ mỗi 64 iteration), một cơ chế tiết kiệm chi phí optimizer hoàn toàn khác, không liên quan tới pruning.}
\end{frame}

% --- Slide 62 ---
\begin{frame}{Quan sát 2: Nút thắt là RAM hệ thống, không phải VRAM}
\begin{itemize}
    \item Đo trên Google Colab T4 trong quá trình huấn luyện:
\end{itemize}
\footnotesize
\begin{table}[]
\centering
\begin{tabular}{lccc}
\toprule
Tài nguyên & Đỉnh sử dụng & Khả dụng & \% sử dụng \\
\midrule
VRAM (GPU) & 1.16 GB & 14.56 GB & \textbf{8\%} \\
RAM hệ thống & 9.59 GB & 12.7 GB & \textbf{91\%} (giới hạn mềm 10.5 GB) \\
\bottomrule
\end{tabular}
\end{table}
\normalsize
\begin{itemize}
    \item RAM hệ thống áp sát giới hạn mềm của pipeline, trong khi VRAM còn dư rất nhiều.
    \item Khuyến nghị cũ "giảm độ phân giải để tiết kiệm VRAM" \textbf{không cần thiết} ở ngân sách này.
    \item \texttt{resolution=1} (độ phân giải gốc) hoàn toàn khả thi và còn loại bỏ luôn độ lệch train/eval resolution.
\end{itemize}
\end{frame}

% --- Slide 63 ---
\begin{frame}{Quan sát 3: Live-score vs reported-score — khác protocol, không phải suy giảm}
\begin{itemize}
    \item Score đo \textbf{trong lúc train} (live) cao hơn hẳn score \textbf{báo cáo cuối} (reported) — nhưng đây không phải mô hình bị tệ đi.
\end{itemize}
\footnotesize
\begin{table}[]
\centering
\begin{tabular}{lccc}
\toprule
Metric & Live (trong train) & Reported (cuối) & Chênh lệch \\
\midrule
PSNR (dB) & 26.23 & 24.46 & $-1.77$ \\
SSIM & 0.8675 & 0.8142 & $-0.0533$ (ít nhất) \\
LPIPS & 0.114 & 0.311 & $+0.197$ \\
\bottomrule
\end{tabular}
\end{table}
\normalsize
\begin{itemize}
    \item Nguyên nhân: (1) LPIPS backbone khác nhau — AlexNet lúc train vs VGG lúc báo cáo; (2) độ phân giải khác nhau — nửa độ phân giải lúc train vs độ phân giải gốc lúc báo cáo.
    \item \alert{Bài học:} không so sánh chéo hai protocol đo khác nhau khi đánh giá tiến bộ giữa các lần chạy.
\end{itemize}
\end{frame}

% --- Slide 64 ---
\begin{frame}{Tổng kết pipeline render 3D Gaussian Splatting}
\centering
\footnotesize
\begin{tikzpicture}[
    node distance=8mm and 6mm,
    box/.style={draw, rounded corners, align=center, minimum height=0.9cm, minimum width=2.6cm, fill=blue!8}
]
\node[box] (gauss) {3D Gaussian\\$(\mu,\Sigma,\alpha,SH)$};
\node[box, right=of gauss] (proj) {Chiếu phối cảnh\\(EWA)};
\node[box, right=of proj] (raster) {Rasterizer khả vi\\theo tile};

\node[box, below=of raster] (blend) {Alpha\\blending};
\node[box, left=of blend] (loss) {Loss\\($L_1$+D-SSIM)};
\node[box, left=of loss] (back) {Backward\\(gradient)};

\node[box, below=of back, fill=orange!12] (adc) {Adaptive\\Density Control};

\draw[-Latex] (gauss) -- (proj);
\draw[-Latex] (proj) -- (raster);
\draw[-Latex] (raster) -- (blend);
\draw[-Latex] (blend) -- (loss);
\draw[-Latex] (loss) -- (back);
\draw[-Latex] (back) -- (adc);
\draw[-Latex] (adc) -- (gauss);
\end{tikzpicture}
\vspace{2mm}

\normalsize
Một chu trình khép kín: render $\rightarrow$ tính loss $\rightarrow$ lan truyền ngược $\rightarrow$ điều chỉnh mật độ Gaussian $\rightarrow$ lặp lại.
\end{frame}

% --- Slide 65 ---
\begin{frame}{Tổng kết cải tiến SADGS \& giới hạn còn lại}
\textbf{Ba đòn bẩy chính và tác động:}
\begin{itemize}
    \item \textbf{Structure-aware densification}: so khớp $\eta$ (screen-space extent / bước sóng texture cục bộ) $\Rightarrow$ split dị hướng đúng trục cần -- \textbf{đang chạy thật} trong \texttt{train.py} qua \texttt{densify\_and\_prune\_structgs}. Riêng bước "giãn Gaussian dưới cỡ bằng giải tích" (\texttt{expand\_undersized\_gs}) \textbf{đã bị comment-out} tại điểm gọi trong \texttt{train.py} -- hàm tồn tại trong \texttt{gaussian\_model.py} nhưng \textbf{không được kích hoạt} ở pipeline mặc định hiện tại.
    \item \textbf{Multiview-consistency} cho cả split lẫn prune (\texttt{split/prune\_ratio\_threshold}) $\Rightarrow$ giảm quyết định "ăn may" theo 1 view nhiễu -- cơ chế này \textbf{đang chạy thật}. Riêng \textbf{multi-view reconstruction consistency score} dùng để loại floater ở \texttt{final\_prune\_structgs} \textbf{không chạy}: lệnh gọi hàm này cũng bị comment-out trong \texttt{train.py} (xem Slide 61) -- đây là thiết kế/dead-code hiện tại, chưa phải cơ chế sống.
    \item \textbf{Compact box} (\texttt{mult}, kiểm soát số tile/splat) + \textbf{Sparse Adam} (\texttt{SparseGaussianAdam}, dùng khi \texttt{optimizer\_type="sparse\_adam"} hoặc mặc định \texttt{"hybrid"}) -- cả hai là cải tiến rasterizer/optimizer nằm sẵn trong \texttt{arguments/\_\_init\_\_.py}, không tồn tại như một module/baseline độc lập có thể chạy riêng trong repo này. Không có "Fused Adam" nào trong mã nguồn -- chỉ có \texttt{SparseGaussianAdam}.
    \item Xem Phần "SADGS (SIGGRAPH 2026)" trong bộ slide này (8 slide) để biết chi tiết công thức $\eta$, split dị hướng, expand giải tích, và pruning đa view.
\end{itemize}
\vspace{2mm}
\textbf{Giới hạn / hướng phát triển:}
\begin{enumerate}
    \item Dữ liệu Gaussian đầu ra chưa được nén (248 B/Gaussian, thô).
    \item \texttt{final\_prune\_structgs} (floater đa-view) là dead code trong \texttt{train.py} hiện tại -- không chạy ở ngân sách nào; ngoài ra ngân sách 7k iteration còn bỏ lỡ luôn mốc prune-opacity đơn giản thứ hai (mặc định \texttt{prune\_iterations=[4000, 8000]}) vì dừng trước 8000 (xem Slide 61).
    \item Chưa có ablation study cô lập riêng từng thành phần ($\eta$ mode, multiview-consistency, expand giải tích) để đo đóng góp thực tế.
    \item Cơ chế \texttt{eta\_compute\_mode="projection"} và \texttt{adaptive\_clone} vẫn ở dạng thử nghiệm, chưa phải cấu hình mặc định khuyến nghị.
    \item Structure-aware densification mới được review theo đúng mã nguồn \texttt{SADGS/scene/gaussian\_model.py} và \texttt{SADGS/utils/freq\_utils.py}, chưa kèm benchmark PSNR/SSIM/LPIPS thật trong bộ slide này.
\end{enumerate}
\end{frame}

% --- Slide 66 ---
\begin{frame}{Bản đồ tài liệu tham khảo}
\footnotesize
\begin{table}[]
\centering
\begin{tabular}{p{4.2cm}p{6.5cm}}
\toprule
\textbf{Tài liệu} & \textbf{Nội dung} \\
\midrule
\texttt{README.md} & Điểm vào tiếng Anh, bảng kết quả đầy đủ \\
\texttt{DOCS/structgs-\allowbreak acceleration-\allowbreak method.md} & Nền tảng toán 3DGS + 3 đòn bẩy SADGS + đo SADGS vs 3DGS \\
\texttt{DOCS/colab-t4-guide.md} & Playbook Colab T4, chống OOM, lộ trình CUDA \\
\texttt{DOCS/BOOK/} & Sách hợp nhất 15 chương, từ toán nền tảng tới cost model \\
\texttt{DOCS/BOOK/*\_figures/} & Biểu đồ minh hoạ công thức cơ chế Structure-Aware Densification (eta, split dị hướng, expand, multiview-consistency) \\
GitHub & \texttt{KietAnhCS/SADGS} \\
\bottomrule
\end{tabular}
\end{table}
\end{frame}

% --- Slide 67 ---
\begin{frame}[c]{}
\centering
{\Huge \textbf{Cảm ơn đã theo dõi!}}

\vspace{6mm}
{\LARGE Hỏi đáp (Q\&A)}

\vspace{10mm}
{\small Digital Twin GS --- SADGS \\ \texttt{github.com/KietAnhCS/SADGS}}
\end{frame}
